% Radial-identity repair / execution appendix for DASHIAtomicPeriodicTable369Formalism.tex
% Repo-native owners:
%   DASHI/Physics/Chemistry/AtomicPeriodicTable369ArchiveHartreeRadialIdentityDefectExact.agda
%   DASHI/Physics/Chemistry/AtomicPeriodicTable369RadialIdentityRepairExecutionExact.agda
%   DASHI/Physics/Chemistry/AtomicPeriodicTable369HydrogenNISTComparisonExact.agda
%   DASHI/Physics/Chemistry/AtomicPeriodicTable369HydrogenReducedMassExact.agda
%   DASHI/Physics/Chemistry/AtomicPeriodicTable369HeliumHartreeConvergenceExact.agda
%   DASHI/Physics/Chemistry/AtomicPeriodicTable369PhysicalValidation.agda

\section{Radial-state identity repair and execution receipt}

The historical \emph{DASHI Atom} Hartree code-shaped object accepted subshell
triples $(n,\ell,\mathrm{occ})$, but passed only $\ell$ into its radial
eigenproblem and always requested the lowest eigenpair in that $\ell$ sector.
Consequently $1s$, $2s$, $3s$, \ldots{} collapsed to the same radial target.
This defect is preserved append-only in
\texttt{AtomicPeriodicTable369ArchiveHartreeRadialIdentityDefectExact.agda}; the
later repair does not rewrite the historical source state.

\subsection{Corrected radial selector}

For an analytically controlled hydrogenic central potential, the radial state
index associated with a bound label $(n,\ell)$ is
\[
  r(n,\ell)=n-\ell-1.
\]
The repaired producer
\[
  \texttt{scripts/atomic\_radial\_identity\_repair.py}
\]
therefore diagonalizes each complete $\ell$ sector far enough to contain every
required radial index and selects eigenpair $r(n,\ell)$. Solving the sector as
one eigenproblem also gives a direct within-$\ell$ orthogonality check rather
than solving the same lowest state repeatedly.

The committed execution receipt is
\[
\texttt{Artifacts/atomic-periodic-table-369/radial-identity-repair-receipt.json}
\]
with producer SHA--256
\[
\texttt{d4df9c2ef1e2292cdf216d7bf5da4aab78189b47cff02dc83bb28d308d240b9e}.
\]
It records Python 3.13.5, NumPy 2.3.5 and SciPy 1.17.0, $4000$ radial grid
points on $r\in(0,80]$ Bohr, and an energy tolerance of $5\times 10^{-4}$
Hartree.

\subsection{Executed hydrogenic regression}

For $Z=1$ and $V_H=0$, the numerical producer checks
\[
  E_n=-\frac{1}{2n^2}\ \text{Hartree}
\]
and radial node count $n-\ell-1$. The retained values are:

\begin{center}
\begin{tabular}{lrrrr}
\toprule
state & radial index & nodes & computed $E$ & analytic $E$\\
\midrule
$1s$ & 0 & 0 & $-0.4999500349782737$ & $-0.5$\\
$2s$ & 1 & 1 & $-0.12499687671803333$ & $-0.125$\\
$3s$ & 2 & 2 & $-0.05555493859367163$ & $-1/18$\\
$2p$ & 0 & 0 & $-0.12500104120998212$ & $-0.125$\\
$3p$ & 1 & 1 & $-0.055556035451670566$ & $-1/18$\\
$3d$ & 0 & 0 & $-0.055555596687603834$ & $-1/18$\\
\bottomrule
\end{tabular}
\end{center}

All energy, node-count and within-$\ell$ orthogonality checks passed. The
largest retained Gram-matrix deviation was below $7\times10^{-16}$.

\subsection{First authoritative H I reference-data diagnostic}

The repaired $1s$ energy is consumed by
\texttt{scripts/atomic\_hydrogen\_nist\_comparison.py}. The committed receipt
uses the 2022 CODATA conversion
\[
  1E_h=27.211386245981\ \mathrm{eV}
\]
and NIST ASD SRD~78
\[
  IE_1(\mathrm H)=13.598434599702\ \mathrm{eV},
\]
with uncertainty $1.2\times10^{-11}$ eV. The finite-grid result is
\[
  IE_1^{\rm grid}=13.604333505485517\ \mathrm{eV},
\]
with residual $+0.005898905783517705$ eV. The exact infinite-nuclear-mass
Coulomb value in the same units is
\[
  IE_1^{\infty M}=13.6056931229905\ \mathrm{eV},
\]
with residual $+0.00725852328850074$ eV. The finite-grid value being closer is
not promoted as better physics: discretization partially cancels omitted model
physics.

\subsection{Leading finite-proton-mass correction}

Using the 2022 CODATA ratio
\[
  \frac{m_e}{m_p}=5.446170214889\times10^{-4},
\qquad
  \frac{\mu}{m_e}=0.9994556794247615,
\]
gives
\[
  IE_1^{\mu}=13.598287264283275\ \mathrm{eV},
\]
with NIST residual
\[
  -0.0001473354187240261\ \mathrm{eV}.
\]
Thus the leading finite-nuclear-mass correction pays most of the large
infinite-mass offset, while relativistic, radiative/QED, higher-order recoil and
proton-size effects remain explicitly unpaid.

\subsection{Executed closed-shell helium continuation}

The first same-object many-electron producer is
\[
\texttt{scripts/atomic\_helium\_hartree\_scf.py}.
\]
It treats neutral helium as a restricted closed-shell $1s^2$ Hartree product.
Each electron sees the Coulomb field of the other electron and the total energy
is
\[
  E_{\rm He}=2\varepsilon_{1s}-J.
\]
At $32000$ radial points it converges in $64$ iterations to
\[
  E_{\rm He}=-2.8604050079530117\ E_h,
\]
with
\[
  IE_1^{\rm grid}(\mathrm{He})=23.412898027945555\ \mathrm{eV},
\]
against NIST ASD
\[
  IE_1(\mathrm{He})=24.587389011\ \mathrm{eV}.
\]
The raw $32000$-point residual is therefore
\[
  -1.1744909830544437\ \mathrm{eV}.
\]

To distinguish discretization from approximation-class error, the convergence
producer
\[
\texttt{scripts/atomic\_helium\_hartree\_convergence.py}
\]
runs the same object at $4000,8000,16000,32000,64000$ grid points. The total
energies monotonically approach the published helium Hartree--Fock/SCF limit
\[
  E_{\rm HF}=-2.861679995612\ E_h
\]
(Mohallem, DOI \texttt{10.1016/0009-2614(92)85634-M}). The two finest grids
produce the first-order Richardson estimate
\[
  E_{\rm Rich}=-2.861682238129302\ E_h,
\]
which differs from the published HF limit by only
\[
  -2.2425173020756972\times10^{-6}\ E_h.
\]
Thus the corrected producer is converging to its intended approximation class.

For a separate correlated benchmark, Pekeris' 1959 helium ground-state value
(doi \texttt{10.1103/PhysRev.115.1216}) is
\[
  E_{\rm Pekeris}=-2.903724375\ E_h.
\]
The HF-to-Pekeris gap is therefore
\[
  0.042044379387999875\ E_h
  \approx 1.1440858469994268\ \mathrm{eV}.
\]
This is the critical defect-routing result: after the radial identity repair and
grid convergence, the dominant surviving helium discrepancy is unpaid
electron correlation/model physics, not recurrence of the archived $n$-erasure
bug.

The execution receipt is
\[
\texttt{Artifacts/atomic-periodic-table-369/helium-hartree-convergence-receipt.json}.
\]

\subsection{Attribution snowball}

The identifier fields remain source-role separated:

\begin{itemize}
\item Schr\"odinger, \emph{Quantisierung als Eigenwertproblem} (1926), DOI
  \texttt{10.1002/andp.19263840404}, is the historical eigenvalue-problem
  source; atom Q9121 and quantum mechanics Q944 are semantic coordinates, with
  Dewey 539.7 and 530.12 as catalogue coordinates.
\item Hartree 1928, DOI \texttt{10.1017/S0305004100011920}, is primary SCF
  provenance. The paper itself reports a helium self-consistent-field
  ionization potential near 24.85 V, explicitly making He a historical SCF
  test object rather than a new DASHI choice.
\item Helium is Wikidata Q560. The current semantic item carries Dewey 553.971;
  this is a catalogue coordinate only and must not be confused with the more
  theory-oriented periodic-table/atomic-physics Dewey route.
\item Hartree--Fock is Q7879841; ionization energy is Q483769 / Gold Book
  I03199. QID identity does not imply model correctness.
\item NIST Atomic Spectra Database SRD~78, DOI \texttt{10.18434/T4W30F}, is the
  downstream authoritative reference-data source. NIST reports helium ground
  shell $1s^2$, ground level $^1S_0$, and $IE_1=24.587389011$ eV with quoted
  uncertainty $2.5\times10^{-8}$ eV.
\item Mohallem 1992, DOI \texttt{10.1016/0009-2614(92)85634-M}, is used only as
  the published helium HF/SCF-limit benchmark. Pekeris 1959, DOI
  \texttt{10.1103/PhysRev.115.1216}, supplies the correlated nonrelativistic
  benchmark. Neither source imports correctness into the DASHI producer.
\item OEIS A320506 records rounded first-ionization energies by atomic number;
  A320507 records atomic numbers sorted by first ionization energy. They are
  useful downstream sequence/navigation coordinates for the future full-$Z$
  residual ladder, not proof or mechanism authority.
\item OEIS A167268, A093907 and A018227 remain upstream structural-selector,
  period-length and closure coordinates. Their role is distinct from A320506
  and A320507, which describe an empirical ionization-energy observable.
\end{itemize}

\subsection{Promotion boundary and next producer}

The paid physical chain is now
\[
\boxed{
\text{radial identity}
\to \text{H execution/reference diagnostic}
\to \text{finite-mass correction}
\to \text{executed He Hartree}
\to \text{He grid convergence to HF limit}.
}
\]
It does \emph{not} imply correlated helium, generic multi-orbital SCF/HF,
correct empirical ground-state configurations, nuclear stability, or full
periodic-table recovery.

The highest-value periodic-table continuation is therefore no longer finer
helium gridding. It is to carry the repaired $(n,\ell)$ state identity into a
multi-orbital SCF/HF producer, beginning with lithium/beryllium-sized occupied
spaces, emit matched neutral/cation total energies, and compare
\[
  IE_1(Z)=E(Z-1)-E(Z)
\]
against NIST ASD while retaining the A320506/A320507 sequence coordinates and
source-specific uncertainties.
